\section{问题一：因素关系与振打峰值}

\subsection{入口条件与操作参数的作用关系}

在参考工况附近分别改变单一变量，得到图\ref{fig:effects}。这些曲线反映式\eqref{eq:fulltwin} 的物理响应，而不是附件直接回归系数。

\begin{figure}[H]
\centering
\includegraphics[width=0.94\textwidth]{fig06_effect_curves.png}
\caption{入口浓度、流量、电压和温度的单因素响应曲线}
\label{fig:effects}
\end{figure}

从式\eqref{eq:fulltwin} 可得以下关系。

\begin{enumerate}
  \item \textbf{入口浓度。} 在去除指数不变时，$C_{\mathrm{out}}$ 与 $C_{\mathrm{in}}$ 近似成正比；同时更高的 $C_{\mathrm{in}}Q$ 增大极板负荷和振打再飞扬，因此实际响应略高于线性。
  \item \textbf{烟气流量。} $Q$ 增大缩短有效停留时间，使 $D\propto Q^{-1}$ 减小，穿透率指数上升。高流量与高入口浓度叠加时风险最强，故采用质量负荷 $C_{\mathrm{in}}Q$ 选择高、低代表工况。
  \item \textbf{电场电压。} $\partial C_{\mathrm{out}}/\partial U_i<0$，且指数模型使边际绝对减排量逐渐递减。按附件中位电压平方归一化后，前两场的基准代理贡献约各占 30\%，后两场约各占 20\%；具体调节顺序由运行电压、功率边际和活跃约束共同决定。
  \item \textbf{温度。} $126.008^\circ\mathrm C$ 是附件均值确定的归一化参考中心，而非由出口数据识别的真实最优温度。钟形先验描述偏离参考温区后的有效迁移能力折减；图\ref{fig:tempsurface} 显示高温与低电压叠加时风险迅速增大，说明高温工况不宜依赖单一末级补偿。
\end{enumerate}

\begin{figure}[H]
\centering
\includegraphics[width=0.88\textwidth]{fig07_temperature_voltage_surface.png}
\caption{温度—电压耦合下的出口浓度响应面}
\label{fig:tempsurface}
\end{figure}

\subsection{振打周期与瞬时峰值}

图\ref{fig:rapping} 给出式\eqref{eq:duststate} 的示意积灰轨迹以及不同周期下的相位叠加峰值。振打周期存在两类相反效应：

\begin{itemize}
  \item 周期延长使单位时间振打次数减少，振打机械能耗下降；但极板尘负荷升高、捕集效率折减，且单次脱落尘饼更大，峰值幅度上升。
  \item 周期缩短降低单次积灰量和单次再飞扬风险，但会增加机械动作频率与振打平均功率；若四场同步振打，即使单场峰值较小，也可能在出口叠加。
\end{itemize}

\begin{figure}[H]
\centering
\includegraphics[width=0.95\textwidth]{fig08_rapping_dynamics.png}
\caption{极板积灰隐状态、振打脉冲与相位叠加示意}
\label{fig:rapping}
\end{figure}

因此，振打周期不宜仅按降低机械能耗的目标取最大值。式\eqref{eq:rapfactor} 保留线性尘量项：周期越长，归一化积灰量和相位平均再飞扬风险越高；周期缩短所增加的机械动作能耗由功率模型中的 $1/T_i$ 项描述。二者共同构成排放风险与机械能耗之间的权衡，式\eqref{eq:dustconstraint} 则限制尘负荷不超过历史高负荷范围。由于缺少实际触发事件，图\ref{fig:rapping} 仅展示归一化相位和相对峰值。现场控制中可将四场触发相位按 $0,\frac14,\frac12,\frac34$ 周期错开，以降低多场峰值叠加风险。

\subsection{问题一结论的可信范围}

附件数据能够支持出口浓度边界特征与能耗关系的分析，但不足以直接辨识各操作量对排放的净效应。因此，入口浓度、流量、电压、温度和振打周期的影响方向由物理结构与文献确定，并通过情景传播和关键参数 $\pm30\%$ 敏感性分析评价其稳定性。所得浓度曲线应理解为具有机理约束的风险外推结果。
